Streaming anomaly detectors tell operators \emph{that} a sensor reading is abnormal and, when interpretable, \emph{which} signal and in \emph{which} direction. They do not tell operators \emph{what to do}: whether to recalibrate, replace, or clean a sensor, or to investigate the process itself. We close this gap with a streaming sensor-fault classifier that attributes each flagged signal to a maintenance-relevant fault type---bias, drift, accuracy loss, or freezing---directly from the conditional residuals of a self-supervised online detector, using O(1) memory and no history buffers. On real battery-depletion faults in the Intel Berkeley Lab deployment, the composed system reaches 90\% real-fault detection at a 7.4\% healthy false-positive rate, and types the resulting faults predominantly as freezing (0.93), the signature of sensors stuck at degraded values. We further characterize the operating envelope on controlled injections: freezing and drift are recovered with recall 1.000 and 0.741, while bias recall is monotone in conditional standard deviation---the interpretable conditional model absorbs constant offsets on strongly coupled signals (recall 0.059) and isolates them on weakly coupled ones (0.941). We show that two original claims do not survive validation: freezing is not a detection blind spot of the base detector (it catches every injected freeze, usually faster than the dedicated test), and adaptation robustness comes from graded change-point suppression rather than residual cancellation. The result is a deployable typing layer with a rigorously mapped---rather than overstated---operating envelope.
